\paragraph{Hint VI.31.E1.} It is an irreducible component of itself.

\paragraph{Hint VI.31.E2.} Use height of $(x)$.

\paragraph{Hint VI.31.E3.} Use the coordinate prime chain.

\paragraph{Hint VI.31.E4.} Use $(0)\subset(x)\subset(x,y)\subset(x,y,z)$.

\paragraph{Hint VI.31.E5.} There is no larger irreducible closed subset inside that component chain.

\paragraph{Hint VI.31.E6.} Local dimension equals height.

\paragraph{Hint VI.31.E7.} The maximal ideal has height two.

\paragraph{Hint VI.31.E8.} Use the generic stalk.

\paragraph{Hint VI.31.E9.} Specialization can only increase prime height.

\paragraph{Hint VI.31.E10.} Use dimension--codimension additivity.

\paragraph{Hint VI.31.E11.} Subtract one.

\paragraph{Hint VI.31.E12.} The dimensions are one and zero.

\paragraph{Hint VI.31.E13.} Apply additivity.

\paragraph{Hint VI.31.E14.} Compare dimensions $2$ and $4$.

\paragraph{Hint VI.31.E15.} It is a coordinate subspace.

\paragraph{Hint VI.31.E16.} Allow a redundant equation.

\paragraph{Hint VI.31.E17.} Use the principal ideal theorem for irreducible components.

\paragraph{Hint VI.31.E18.} It is an irreducible component.

\paragraph{Hint VI.31.E19.} Both terms vanish.

\paragraph{Hint VI.31.E20.} Use nested additivity for varieties.

\paragraph{Hint VI.31.E21.} Compare underlying topological spaces.

\paragraph{Hint VI.31.E22.} Take the infimum over component codimensions.

\paragraph{Hint VI.31.E23.} Translate zero chain length into maximality.

\paragraph{Hint VI.31.E24.} Recall the chapter boundary.

\paragraph{Hint VI.31.E25.} VI/30 consumed (d), VI/31 consumes (c).

\paragraph{Hint VI.31.E26.} One is absolute; the other is relative to a chosen closed locus.